\section{System and Protocol Model}
\label{sec:system}

\subsection{Entities and Controlled Resource}

The system contains publishers, an MQTT~5 broker, a co-located application-state controller, and LLM subscribers. The controller runs as a separate Node.js process on the broker host and manages \emph{gateway/application-state capacity}; Mosquitto's native in-flight store remains unchanged. Co-location removes an avoidable wide-area hop but is not required by the protocol.

For request $i$, the controller stores
\begin{equation}
 r_i=(id_i,c_i,t_i,D_i,x_i,z_i),
\end{equation}
where $id_i$ is the correlation identifier, $c_i\in\{g,c\}$ is the assigned GPU or CPU-only class, $t_i$ is admission time, $D_i$ is the application deadline, $x_i$ is the inference configuration, and $z_i$ is state. All records occupy one unit in this implementation; unequal byte sizes and multirate loss control are left to future work. With capacity $n$, an arrival that sees $n$ records is rejected immediately and receives an application response with status \texttt{blocked}.

Let $S_{i}$ be elapsed time from admission until the controller observes the correlated completion. It includes broker-to-subscriber delivery, subscriber queueing, inference, and result publication. The record holding time is
\begin{equation}
 H_i=\min(S_i,D_i).
 \label{eq:holding}
\end{equation}
If $S_i>D_i$, the record is released on expiry; a later result is logged as late and does not recreate state. Correct, on-time goodput counts completions that arrive before expiry, parse as valid JSON, and contain the correct task label.

\subsection{Protocol Sequence}

Figure~\ref{fig:sequence} shows the complete path. A publisher first subscribes to its response topic and sends a QoS~1 PUBLISH request containing Response Topic and Correlation Data. Mosquitto returns PUBACK according to MQTT rules. The controller independently receives the request event and either rejects it or creates $r_i$. On admission, it publishes a job to the selected subscriber-class topic with three MQTT~5 properties: Correlation Data, Message Expiry Interval, and User Properties encoding the application deadline, configuration, and class.

The subscriber processes the job and publishes the result with the same Correlation Data. The controller correlates the result, releases $r_i$, and publishes a completion response. If the timer fires first, it releases $r_i$ and publishes an expiry response. The publisher-facing response is therefore an application message, never PUBACK/PUBREC/PUBCOMP.

\begin{figure*}[t]
  \centering
  \includegraphics[width=0.98\textwidth]{fig5_mqtt_sequence.png}
  \caption{MQTT--LLM application-state control workflow. The MQTT QoS acknowledgment is independent of application-level completion. The controller maintains bounded request state and releases it upon correlated LLM completion or deadline expiry.}
  \label{fig:sequence}
\end{figure*}

Message Expiry Interval remains useful but has a narrower role. Per MQTT~5, the broker removes an expired message if onward delivery has not begun~\cite{oasis2019mqtt}. It does not keep a delivered request resident while an LLM executes. Equation~\eqref{eq:holding} is implemented by the controller's timer and completion correlation.

\subsection{Traffic and Subscriber Classes}

During a short control epoch $\Delta$, the feedforward model treats admitted attempts as a locally stationary Poisson process with rate $\lambda$. This approximation is intended for many independent, low-rate sources with randomized phases and no dominant source. Synchronized publishers, alarms, retries, and small source populations fall outside this regime. The implementation records the one-second index of dispersion
\begin{equation}
 I=\frac{\operatorname{Var}[N(1\ \mathrm{s})]}{\mathbb{E}[N(1\ \mathrm{s})]},
\end{equation}
and the evaluation treats large deviations from $I=1$ as evidence of mismatch rather than discarding the run.

The two subscriber classes execute the same model and task but use different Ollama endpoints: one may use the RTX~3060 GPU, while the other is explicitly forced to CPU-only execution on the same host. The principal live matrix uses Qwen2.5--3B, and the drift condition changes both endpoints to Qwen2.5--7B halfway through a run. This design isolates execution-class heterogeneity without falsely claiming two physical edge devices. Class preference is drawn independently for each arrival with CPU fraction $p_c$; the guarded policy may reroute a request if its preferred queue is predicted to miss a class-specific feasibility floor.

\subsection{Objective and Metrics}

The controller operates under three competing outcomes:
\begin{align}
 B_{\mathrm{app}} &= \frac{N_{\mathrm{blocked}}}{N_{\mathrm{published}}},\\
 E_{\mathrm{adm}} &= \frac{N_{\mathrm{expired}}}{N_{\mathrm{admitted}}},\\
 G_Q &= \frac{N_{\mathrm{correct,on\mbox{-}time}}}{T_{\mathrm{run}}}.
\end{align}
Reducing $D$ may reduce $B_{\mathrm{app}}$ while increasing $E_{\mathrm{adm}}$; increasing it may do the reverse. Hence blocking alone is not a sufficient objective. We report all three metrics, strict on-time accuracy, MQTT PUBACK latency, end-to-end completion latency, occupancy, rerouting, and process memory. The design goal is high $G_Q$ with explicit visibility into whether capacity is protected by successful completion, expiry, or admission shedding.
